\documentclass[
a4paper,
amsfonts,
amssymb,
amsmath,
reprint,
twoside
]{revtex4-1}

% =========================================================
% CONFIGURACIÓN IDIOMA Y CODIFICACIÓN
% =========================================================

\usepackage[spanish]{babel}
\usepackage[utf8]{inputenc}
\usepackage[T1]{fontenc}

% =========================================================
% PAQUETES MATEMÁTICOS
% =========================================================

\usepackage{amsmath}
\usepackage{amssymb}
\usepackage{bm}

% =========================================================
% FIGURAS Y TABLAS
% =========================================================

\usepackage{graphicx}
\usepackage{float}
\usepackage{booktabs}
\usepackage{multirow}
\usepackage{array}

% =========================================================
% HIPERVÍNCULOS
% =========================================================
\usepackage[
pdftex,
pdftitle={Taller de reconocimiento de materiales},
pdfauthor={Nombres de los integrantes}
]{hyperref}
% =========================================================
% REDUCCIÓN ESPACIADO ECUACIONES
% =========================================================

\AtBeginDocument{

\setlength{\abovedisplayskip}{6pt}
\setlength{\belowdisplayskip}{6pt}

\setlength{\abovedisplayshortskip}{4pt}
\setlength{\belowdisplayshortskip}{4pt}

}

% =========================================================
% REDUCCIÓN ESPACIADO LISTAS
% =========================================================

\usepackage{enumitem}

\setlist[itemize]{
noitemsep,
topsep=0pt,
parsep=0pt,
partopsep=0pt
}

% =========================================================
% FECHA
% =========================================================

\makeatletter
\def\Dated@name{}
\makeatother

% =========================================================
% ESTILO BIBLIOGRÁFICO
% =========================================================

\bibliographystyle{apsrev4-1}

% =========================================================
% DOCUMENTO
% =========================================================

\begin{document}

% =========================================================
% TÍTULO
% =========================================================

\title{
Simulación computacional de un proceso de guerra de guerrillas
}

% =========================================================
% AUTORES
% =========================================================

\author{
Juan Esteban Gómez Gómez, Jose Luis Martinez Diaz, William David Rincón Verano, Serena Triviño Sáenz
}

\affiliation{
Universidad Nacional de Colombia
}

\date{30 de octubre de 2026}

% =========================================================
% RESUMEN
% =========================================================

%------------------------------------------------
\begin{abstract}

En este trabajo se estudia mediante simulación computacional un
modelo estocástico de fragmentación y fusión aplicado a un proceso
de guerra de guerrillas. El sistema se representa mediante unidades
de ataque que interactúan de acuerdo con reglas probabilísticas,
conservando una fuerza total de $N=30\,000$. Se desarrollaron dos
implementaciones independientes, en C++ y Python, para obtener
las distribuciones de tamaños y analizar su comportamiento frente
a una ley de potencia de referencia, caracterizada por el exponente
teórico $\alpha=2.5$. Asimismo, se examinó la influencia del
parámetro de fragmentación $\nu$ para cuatro valores entre
$0.005$ y $0.1$. Los resultados muestran distribuciones decrecientes
compatibles con un comportamiento de ley de potencia y exponentes
estimados entre $2.514$ y $2.634$ en Python, y entre $2.5686$
y $2.5953$ en C++. La comparación entre ambas implementaciones
presentó diferencias relativas inferiores al $2.2\%$. Finalmente,
se verificó la conservación de la fuerza total en Python y se
compararon los tiempos de ejecución, observándose un menor tiempo
registrado para C++ bajo los procedimientos computacionales
utilizados.\\

Palabras clave: fragmentación, fusión,
simulación estocástica, guerra de guerrillas, ley de potencia.
\end{abstract}

\maketitle

%------------------------------------------------

%------------------------------------------------
\section{Introducción}

Los sistemas constituidos por múltiples unidades que interactúan
mediante procesos de fragmentación y fusión pueden presentar
distribuciones estadísticas de tamaño que siguen leyes de potencia.
El estudio de estos sistemas permite analizar cómo reglas locales
y probabilísticas generan comportamientos colectivos a escala
global. Johnson et al.~\cite{johnson2005} estudiaron regularidades
estadísticas en conflictos armados y propusieron un marco para
analizar la distribución de la intensidad de los eventos violentos.
Posteriormente, Johnson et al.~\cite{johnson2006} examinaron patrones
estadísticos comunes en guerras y terrorismo, proporcionando un
contexto para el estudio de modelos de fragmentación y fusión.

En este trabajo se estudia un modelo estocástico en el que la fuerza
total del sistema se distribuye entre diferentes unidades de ataque.
Estas unidades pueden fragmentarse en componentes de tamaño unitario
o fusionarse con otras unidades, de acuerdo con reglas probabilísticas.
La evolución del sistema se simula computacionalmente siguiendo el
modelo descrito en los artículos de referencia.

El objetivo es obtener y analizar la distribución de tamaños de las
unidades, contrastarla con la ley de potencia teórica
$n_s \propto s^{-\alpha}$, cuyo exponente de referencia es
$\alpha=2.5$, y estudiar cómo varía la distribución al modificar
el parámetro de fragmentación $\nu$. Para ello, se comparan dos
implementaciones independientes, desarrolladas en C++ y Python,
considerando tanto los resultados estadísticos como sus tiempos
de ejecución.


%------------------------------------------------

%------------------------------------------------
\section{Objetivos}

\subsection{Objetivo general}

Implementar y analizar computacionalmente un modelo estocástico
de fragmentación y fusión aplicado a un proceso de guerra de
guerrillas, con el propósito de estudiar la distribución de tamaños
de las unidades, contrastarla con la ley de potencia teórica de
exponente $\alpha=2.5$ y comparar los resultados y el desempeño
computacional de dos implementaciones desarrolladas en C++ y Python.

\subsection{Objetivos específicos}

\begin{itemize}

    \item Implementar las reglas probabilísticas de selección,
    fragmentación y fusión mediante algoritmos desarrollados
    en C++ y Python.

    \item Verificar la conservación de la fuerza total del sistema
    durante la evolución de la simulación.

    \item Obtener y analizar la distribución estadística del número
    de unidades en función de su tamaño, y contrastarla con el
    comportamiento teórico esperado.

    \item Estimar el exponente $\alpha$ mediante un ajuste lineal
    en escala logarítmica y compararlo con el valor teórico
    de referencia.

    \item Examinar la variación de las distribuciones de tamaños
    y del exponente estimado al modificar el parámetro de
    fragmentación $\nu$.

    \item Comparar los resultados estadísticos y los tiempos de
    ejecución de las implementaciones en C++ y Python,
    considerando las condiciones de simulación utilizadas.

\end{itemize}

%------------------------------------------------

%------------------------------------------------
\section{Modelo matemático}

\subsection{Representación del sistema}

El sistema está constituido por un conjunto de unidades, cada una
caracterizada por un tamaño o fuerza de ataque $s_i$. La fuerza total
del sistema se conserva durante la evolución y se expresa como

\begin{equation}
    \sum_{i=1}^{M}s_i=N,
    \label{eq:conservacion}
\end{equation}

donde $N$ representa la fuerza total del sistema y $M$ es el número
de unidades presentes en un instante determinado.

Inicialmente, el sistema se compone de $N$ unidades de tamaño
unitario, de modo que $s_i=1$ para todas ellas. Las interacciones
modifican el número de unidades y sus tamaños, sin alterar la suma
total de sus tamaños.

\subsection{Selección de unidades}

La selección de una unidad depende de su tamaño. La probabilidad
de seleccionar la unidad $i$ se define como

\begin{equation}
    P_i=\frac{s_i}{N}.
    \label{eq:seleccion}
\end{equation}

Por tanto, las unidades de mayor tamaño tienen una probabilidad
proporcionalmente mayor de ser seleccionadas. Esta regla introduce
una dependencia entre el tamaño de cada unidad y su participación
en las interacciones del sistema.

\subsection{Fragmentación}

Con probabilidad $\nu$, la unidad seleccionada se fragmenta en
$s_i$ unidades de tamaño unitario. Si la unidad original tiene
tamaño $s_i$, la fragmentación conserva la fuerza total, puesto que

\begin{equation}
    s_i=\underbrace{1+1+\cdots+1}_{s_i\text{ unidades}}.
\end{equation}

Este proceso incrementa el número de unidades y distribuye la
fuerza entre componentes de menor tamaño.

\subsection{Fusión}

Con probabilidad $1-\nu$, se selecciona otra unidad de tamaño $s_j$
y ambas se fusionan en una unidad cuyo tamaño es

\begin{equation}
    s_{\mathrm{nueva}}=s_i+s_j.
    \label{eq:fusion}
\end{equation}

La fusión reduce el número de unidades y concentra la fuerza en
componentes de mayor tamaño. La selección de la segunda unidad
se realiza de acuerdo con la regla probabilística implementada
en el algoritmo.

\subsection{Distribución de tamaños y exponente}

Sea $n_s$ el número de unidades de tamaño $s$. La distribución de
tamaños permite caracterizar el estado estadístico del sistema.
Como referencia teórica, se considera una distribución de ley de
potencia de la forma

\begin{equation}
    n_s\propto s^{-\alpha},
    \label{eq:ley_potencia}
\end{equation}

donde $\alpha$ es el exponente que caracteriza el decrecimiento
de la distribución.

Al tomar logaritmos en ambos lados de la ecuación
\eqref{eq:ley_potencia}, se obtiene una relación lineal:

\begin{equation}
    \log(n_s)=C-\alpha\log(s),
    \label{eq:ajuste_lineal}
\end{equation}

donde $C$ es una constante. Por tanto, el exponente se estima
como el negativo de la pendiente del ajuste lineal entre
$\log(s)$ y $\log(n_s)$.

Para el análisis se utiliza $\alpha=2.5$ como valor teórico de
referencia. La comparación con este valor se realiza mediante
las distribuciones obtenidas y los exponentes estimados
numéricamente.


%------------------------------------------------
\section{Algoritmo e implementación computacional}

\subsection{Procedimiento de simulación}

El modelo se implementó de manera independiente en C++ y Python.
Ambas implementaciones representan los procesos de selección,
fragmentación y fusión de las unidades, con el propósito de
estudiar la evolución del sistema y obtener la distribución de
tamaños de las unidades de ataque.

El procedimiento general consiste en inicializar el sistema con
$N$ unidades de tamaño uno, seleccionar unidades de acuerdo con
la regla probabilística del modelo y decidir si ocurre
fragmentación o fusión mediante el parámetro $\nu$. Después de
cada interacción se actualiza el estado del sistema y se repite
el proceso durante el número de pasos establecido. Finalmente,
se obtiene la distribución de tamaños y se estima el exponente
de escala.

Las implementaciones difieren en la forma de procesar los datos:
Python analiza la distribución del estado final de la simulación,
mientras que C++ acumula muestras durante la ejecución. Esta
diferencia debe tenerse en cuenta al comparar las distribuciones
y los exponentes estimados.

\subsection{Organización y gestión del trabajo colaborativo}

Para coordinar el desarrollo del proyecto se utilizó la plataforma
\textit{Jira}, con el propósito de organizar las actividades,
asignar responsabilidades y realizar seguimiento al avance
del equipo de trabajo. Mediante un tablero colaborativo se
distribuyeron las tareas relacionadas con la revisión del modelo
matemático, la implementación computacional, el análisis de
resultados y la elaboración del informe.

Esta herramienta permitió visualizar las actividades pendientes,
en ejecución y completadas, facilitando la coordinación entre
los integrantes y el seguimiento de los compromisos establecidos.
Como evidencia del proceso de organización se presenta el tablero
utilizado durante el desarrollo del proyecto.

El espacio de trabajo puede consultarse en el siguiente enlace:
\url{https://unal-team-bomb75.atlassian.net/jira/software/projects/KAN/list?jql=project%20%3D%20KAN%20ORDER%20BY%20cf%5B10019%5D%20ASC}.

\begin{figure}[htbp]
    \centering
    \includegraphics[width=1.0\columnwidth]{Figuras/jira.png}
    \caption{Tablero de Jira utilizado para la asignación,
    organización y seguimiento de las actividades del proyecto.}
    \label{fig:tablero_jira}
\end{figure}

\subsection{Parámetros de simulación}

En las simulaciones finales se utilizó una fuerza total
$N=30\,000$ y un parámetro de fragmentación $\nu=0.01$ para
el análisis principal. También se estudiaron cuatro valores
del parámetro de fragmentación:
$\nu\in\{0.005,\,0.01,\,0.05,\,0.1\}$.

La Tabla~\ref{tab:parametros} resume los parámetros principales
de las dos implementaciones.

\begin{table}[htbp]
    \centering
    \caption{Parámetros de las simulaciones principales.}
    \label{tab:parametros}
    \begin{tabular}{lll}
        \toprule
        Parámetro & Símbolo & Valor \\
        \midrule
        Fuerza total & $N$ & $30\,000$ \\
        Fragmentación & $\nu$ & $0.01$ \\
        Pasos de Python & $T_{\mathrm{Py}}$ & $9\,000\,000$ \\
        Pasos de C++ & $T_{\mathrm{C++}}$ & $900\,000\,000$ \\
        Semilla inicial de C++ & -- & $6767$ \\
        Valores de $\nu$ & -- &
        $0.005,\,0.01,\,0.05,\,0.1$ \\
        \bottomrule
    \end{tabular}
\end{table}

En Python, la simulación principal se ejecutó durante 300 barridos,
equivalentes a $9\,000\,000$ eventos para $N=30\,000$. En C++,
la ejecución principal utilizó $900\,000\,000$ pasos. Para el
barrido de $\nu$, se consideraron los cuatro valores indicados.
El número de pasos y el procedimiento de acumulación de muestras
no son iguales en ambas implementaciones, por lo que estas
diferencias se consideran al interpretar los resultados.

\subsection{Estimación del exponente}

El exponente $\alpha$ se estimó mediante un ajuste lineal de
$\log_{10}(n_s)$ frente a $\log_{10}(s)$, considerando los valores
positivos de $n_s$ en el intervalo

\begin{equation}
    2\leq s\leq 30.
\end{equation}

Si $m$ representa la pendiente obtenida mediante regresión lineal,
el exponente estimado se calcula como

\begin{equation}
    \widehat{\alpha}=-m.
    \label{eq:alpha_estimado}
\end{equation}

Este procedimiento proporciona una estimación exploratoria del
exponente. Su valor depende de la distribución obtenida, del
intervalo de ajuste y del procedimiento utilizado para acumular
los datos. Por ello, los resultados se comparan con el valor
teórico de referencia $\alpha=2.5$ sin asumir una coincidencia
exacta.


%------------------------------------------------
\section{Resultados}

\subsection{Distribución de tamaños}

Las simulaciones permiten obtener la distribución del número de
unidades $n_s$ en función de su tamaño $s$. Para el caso principal,
se utilizó una fuerza total $N=30\,000$ y un parámetro de
fragmentación $\nu=0.01$. Las distribuciones obtenidas mediante
las implementaciones en C++ y Python se presentan en las
Figuras~\ref{fig:distribucion_cpp} y~\ref{fig:distribucion_python},
respectivamente.

\begin{figure}[htbp]
    \centering
    \includegraphics[width=1.0\columnwidth]{Figuras/c_1.png}
    \caption{Distribución del número de unidades en función de su
    tamaño, obtenida mediante la implementación en C++ para
    $N=30\,000$ y $\nu=0.01$.}
    \label{fig:distribucion_cpp}
\end{figure}

\begin{figure}[htbp]
    \centering
    \includegraphics[width=1.0\columnwidth]{Figuras/p_1.png}
    \caption{Distribución del número de unidades en función de su
    tamaño, obtenida mediante la implementación en Python para
    $N=30\,000$ y $\nu=0.01$.}
    \label{fig:distribucion_python}
\end{figure}

En ambas implementaciones se observa que el número de unidades
disminuye al aumentar su tamaño, comportamiento compatible con
una distribución de tipo ley de potencia. Sin embargo, las
distribuciones presentan diferencias, por lo que la comparación
cuantitativa se realiza mediante los exponentes estimados.

\subsection{Efecto del parámetro de fragmentación}

Para examinar la influencia de la fragmentación sobre la
distribución de tamaños, se consideraron cuatro valores del
parámetro: $\nu=0.005$, $0.01$, $0.05$ y $0.1$. Los resultados
obtenidos mediante las implementaciones en C++ y Python se
presentan en las Figuras~\ref{fig:distribuciones_nu_cpp}
y~\ref{fig:distribuciones_nu_python}, respectivamente.

\begin{figure}[htbp]
    \centering
    \includegraphics[width=1.0\columnwidth]{Figuras/c_2.png}
    \caption{Distribuciones del número de unidades en función de
    su tamaño, obtenidas mediante C++ para los valores del
    parámetro de fragmentación
    $\nu\in\{0.005,\,0.01,\,0.05,\,0.1\}$.}
    \label{fig:distribuciones_nu_cpp}
\end{figure}

\begin{figure}[htbp]
    \centering
    \includegraphics[width=1.0\columnwidth]{Figuras/p_2.png}
    \caption{Distribuciones del número de unidades en función de
    su tamaño, obtenidas mediante Python para los valores del
    parámetro de fragmentación
    $\nu\in\{0.005,\,0.01,\,0.05,\,0.1\}$.}
    \label{fig:distribuciones_nu_python}
\end{figure}

Los resultados permiten comparar cómo cambia la distribución
al modificar la probabilidad de fragmentación. La magnitud de
estos cambios se evalúa también mediante el exponente
$\alpha$ estimado para cada valor de $\nu$.

\subsection{Estimación y comparación del exponente}

El exponente se estimó mediante un ajuste lineal en escala
logarítmica, utilizando los valores positivos de $n_s$ en el
intervalo $2\leq s\leq 30$. La Tabla~\ref{tab:exponentes}
presenta los resultados de ambas implementaciones y la diferencia
relativa entre sus estimaciones.

\begin{table}[htbp]
    \centering
    \caption{Exponentes estimados mediante C++ y Python y
    diferencia relativa entre los resultados.}
    \label{tab:exponentes}
    \begin{tabular}{cccc}
        \toprule
        $\nu$ & $\alpha_{\mathrm{C++}}$ &
        $\alpha_{\mathrm{Python}}$ & Diferencia (\%) \\
        \midrule
        0.005 & 2.5953 & 2.605 & 0.37 \\
        0.010 & 2.5762 & 2.586 & 0.38 \\
        0.050 & 2.5686 & 2.514 & 2.17 \\
        0.100 & 2.5895 & 2.634 & 1.69 \\
        \bottomrule
    \end{tabular}
\end{table}

La diferencia relativa se calculó respecto al resultado de
Python mediante

\begin{equation}
    \Delta_{\alpha}=
    \frac{\left|\alpha_{\mathrm{C++}}-
    \alpha_{\mathrm{Python}}\right|}
    {\alpha_{\mathrm{Python}}}\times 100\%.
\end{equation}

Las diferencias entre las estimaciones se encuentran entre
aproximadamente $0.37\%$ y $2.17\%$. El menor desacuerdo se
presenta para $\nu=0.005$ y $\nu=0.01$, mientras que el mayor
corresponde a $\nu=0.05$. Como referencia, el exponente teórico
considerado en este trabajo es $\alpha=2.5$. Los valores
numéricos obtenidos se encuentran por encima de dicha referencia,
aunque su cercanía varía según el parámetro y la implementación.

\subsection{Desempeño computacional y conservación}

La ejecución principal en Python utilizó $9\,000\,000$ eventos
y registró un tiempo de $28\,406.1$ segundos. Al finalizar,
el sistema contenía $15\,729$ unidades y la suma de sus tamaños
era $30\,000$, verificando la conservación de la fuerza total:

\begin{equation}
    \sum_{i=1}^{M}s_i=N=30\,000.
\end{equation}

Por su parte, la implementación en C++ registró un tiempo de
ejecución de $188.419$ segundos para $900\,000\,000$ pasos,
con $120\,000$ muestras acumuladas. Estos tiempos corresponden
a los procedimientos utilizados en cada implementación. Dado
que Python analiza el estado final, mientras que C++ acumula
muestras durante la ejecución, la comparación de rendimiento
debe interpretarse teniendo en cuenta esta diferencia
metodológica.

%------------------------------------------------
\section{Discusión}

Los resultados obtenidos muestran que las distribuciones de
tamaños presentan un comportamiento decreciente, en el que las
unidades de mayor tamaño son menos numerosas que las de menor
tamaño. Este comportamiento es compatible con una distribución
de ley de potencia, como la propuesta en el modelo de referencia.
Sin embargo, los exponentes estimados no coinciden exactamente
con el valor teórico $\alpha=2.5$, lo que evidencia diferencias
entre el comportamiento numérico y la predicción teórica.

La comparación entre las implementaciones en C++ y Python muestra
una concordancia general en las estimaciones del exponente.
Las diferencias relativas calculadas se encuentran entre
$0.37\%$ y $2.17\%$, lo que indica que ambos procedimientos
producen resultados cuantitativamente cercanos. Las discrepancias
pueden relacionarse con las diferencias metodológicas entre las
implementaciones: Python obtiene la distribución a partir del
estado final de la simulación, mientras que C++ acumula muestras
durante su ejecución. Por tanto, los resultados no corresponden
a procedimientos de muestreo idénticos.

En cuanto al parámetro de fragmentación $\nu$, los exponentes
estimados presentan variaciones moderadas para los cuatro valores
estudiados. En Python, el exponente disminuye de $2.605$ a
$2.514$ al pasar de $\nu=0.005$ a $\nu=0.05$, y posteriormente
aumenta hasta $2.634$ para $\nu=0.1$. En C++, las estimaciones
se mantienen en un intervalo más estrecho, entre $2.5686$ y
$2.5953$. Estos resultados muestran que, dentro del conjunto
de parámetros analizado, no se observa una variación estrictamente
monótona del exponente con respecto a $\nu$.

Por otra parte, la simulación en Python permitió verificar la
conservación de la fuerza total del sistema, cuya suma permaneció
en $N=30\,000$ al finalizar el procedimiento. Esta comprobación
respalda la correcta implementación de las operaciones de
fragmentación y fusión desde el punto de vista de la conservación.

Finalmente, los tiempos de ejecución registrados muestran una
diferencia importante entre las implementaciones. C++ completó
su procedimiento en $188.419$ segundos, mientras que Python
registró $28\,406.1$ segundos. Aunque estos resultados evidencian
un menor tiempo de ejecución para C++, la comparación debe
considerarse con cautela, debido a que ambas implementaciones
utilizaron diferentes cantidades de pasos y distintos
procedimientos de acumulación de muestras.


%------------------------------------------------
\section{Repositorio y disponibilidad del código fuente}

Con el propósito de organizar y facilitar la consulta de los
recursos computacionales desarrollados, se creó un repositorio
público en GitHub para el proyecto. Este contiene las
implementaciones del modelo en C++ y Python, los cuadernos
Jupyter, las figuras generadas durante las simulaciones,
los resultados numéricos y los archivos de configuración
y documentación.

El repositorio se encuentra disponible en el siguiente enlace:

\begin{center}
    \url{https://github.com/jomartinezdi-UNAL/Proyecto_G_Guerrillas}
\end{center}

La estructura del repositorio permite identificar por separado
el código fuente, los notebooks de ejecución, las figuras,
los resultados y los recursos necesarios para consultar
y reproducir el procedimiento computacional.
%------------------------------------------------


%------------------------------------------------
\section{Conclusiones}

\begin{itemize}

    \item Se implementó computacionalmente un modelo estocástico
    de fragmentación y fusión en dos lenguajes de programación,
    C++ y Python, permitiendo estudiar la evolución de las
    unidades y sus distribuciones de tamaños.

    \item Las distribuciones obtenidas presentan un comportamiento
    decreciente compatible con una ley de potencia. Los exponentes
    estimados se encuentran por encima del valor teórico de
    referencia $\alpha=2.5$, con valores entre $2.514$ y $2.634$
    para Python, y entre $2.5686$ y $2.5953$ para C++.

    \item La comparación entre las implementaciones muestra una
    concordancia cuantitativa en la estimación del exponente,
    con diferencias relativas inferiores al $2.2\%$ para los
    cuatro valores de $\nu$ estudiados. Las discrepancias
    observadas deben interpretarse considerando que los
    procedimientos de acumulación de datos son diferentes.

    \item La variación del parámetro de fragmentación
    $\nu\in\{0.005,\,0.01,\,0.05,\,0.1\}$ produjo cambios
    moderados en los exponentes estimados. No se identificó
    una relación monótona entre dicho parámetro y el exponente
    $\alpha$ dentro del intervalo analizado.

    \item La implementación en Python permitió verificar la
    conservación de la fuerza total del sistema, manteniéndose
    la suma de los tamaños igual a $N=30\,000$ al finalizar
    la simulación.

    \item Finalmente, C++ presentó un menor tiempo de ejecución
    registrado que Python. No obstante, debido a las diferencias
    en el número de pasos y en los procedimientos de muestreo,
    esta comparación no constituye una evaluación controlada
    del rendimiento de ambos lenguajes.

\end{itemize}

%------------------------------------------------
\begin{thebibliography}{9}

\bibitem{johnson2005}
N. Johnson et al.,
\textit{From old wars to new wars and global terrorism},
arXiv:physics/0506213, 2005.

\bibitem{johnson2006}
N. Johnson et al.,
\textit{Universal patterns underlying ongoing wars and terrorism},
arXiv:physics/0605035, 2006.

\end{thebibliography}

\end{document}
